\begin{theorem}[Isometry of the normalized averaging site]%
    \label{thm:peps_averaging_site_isometry}
    \lean{TNLean.PEPS.isGIsometric_averagingSite}
    \leanok
    \uses{def:peps_regular_averaging_site}
    Let $U$ be a matrix representation of a finite group and let $P$
    be the normalized average of its four-leg action. The corresponding
    site map is $G$-isometric, with factor one on the invariant virtual
    subspace: $\langle Px,Py\rangle=\langle x,y\rangle$ for invariant
    $x,y$. This is an auxiliary projector consequence of the isometry
    condition in \cite[Definition~6.1]{Schuch2010PEPS}.
\end{theorem}
\begin{proof}\leanok
    The site map is the averaging map. It is invariant, injective on
    invariant vectors, and acts as the identity on that subspace.
\end{proof}

\begin{theorem}[Exact norm of the actual regular torus state]%
    \label{thm:peps_regular_averaging_torus_norm}
    \lean{TNLean.PEPS.torusBondNetwork_regular_averagingSite_norm,
        TNLean.PEPS.torusBondNetwork_regular_averagingSite_ne_zero}
    \leanok
    \uses{def:peps_regular_averaging_site, lem:peps_isometric_adjoint,
        thm:peps_regular_torus_closure_overlap}
    Let $G$ be finite, let the torus have positive periods, and let $N$
    be its number of sites. For the actual torus contraction
    $\Psi_{\mathrm{reg}}$ of normalized regular averaging sites with
    identity bond operators,
    \begin{align}
        \langle\Psi_{\mathrm{reg}},\Psi_{\mathrm{reg}}\rangle
            &=|G|^{N+1},\\
        \Psi_{\mathrm{reg}}&\ne0.
    \end{align}
    This is an auxiliary normalization consequence for the regular state
    in \cite[Section~7]{Schuch2010PEPS}, with unnormalized virtual bonds.
\end{theorem}
\begin{proof}\leanok
    The regular averaging projector is self-adjoint and idempotent,
    so its actual site Gram kernel has factor one.
    In the regular torus Gram contraction with trivial closures, every
    group element is a simultaneous closure intertwiner. The overlap is
    therefore $|G|^N|G|$. The group order is positive, proving nonvanishing.
\end{proof}

\begin{theorem}[Normalization under the physical state isometry]%
    \label{thm:peps_weighted_torus_norm}
    \lean{TNLean.PEPS.norm_ne_zero_of_isometry_to_regularTorus,
        TNLean.PEPS.torusBondNetwork_blockFourthRootWeight_norm_ne_zero}
    \leanok
    \uses{thm:peps_regular_averaging_torus_norm,
        thm:peps_torus_semiregular_isometric_extension,
        lem:peps_torus_physical_bond_regrouping_overlap}
    Every norm-preserving physical map carrying a vector $\psi$ to the
    actual regrouped regular torus state gives
    $\langle\psi,\psi\rangle=|G|^{N+1}$ and $\psi\ne0$.
    In particular, choose positive-dimensional blocks $D_i$ and an
    isometric Fourier intertwiner with the regular representation.
    The actual fourth-root-weighted block state $\Psi_{U,W}$ then satisfies
    \begin{align}
        \langle\Psi_{U,W},\Psi_{U,W}\rangle
            &=|G|^{N+1},\\
        \Psi_{U,W}&\ne0.
    \end{align}
    Both conclusions concern finite tori with positive periods.
\end{theorem}
\begin{proof}\leanok
    Move the norm-preserving map across the overlap and substitute the
    actual regular-state equality. Physical coordinate regrouping preserves
    overlaps, giving the exact regular norm. Its positivity implies that
    the original vector is nonzero. For the weighted block construction,
    use the already derived ambient physical state isometry.
\end{proof}
